\clearpage
\begin{summarybox}
\paragraph{为什么需要这个结论？}
\textbf{把不好手算的指数、对数、三角函数值，换成容易计算的分式；
再用导数和不等式决定估算的可信程度。}
它是\textbf{接近 $0$ 时的计算工具}，不是新的函数恒等式。

\paragraph{五式速记}
\[
\begin{aligned}
e^x&\approx\frac{x^2+6x+12}{x^2-6x+12},&
\ln(1+x)&\approx\frac{3x^2+6x}{x^2+6x+6},\\
\sin x&\approx\frac{6x}{x^2+6},&
\cos x&\approx\frac{12-5x^2}{12+x^2},\\
\tan x&\approx\frac{3x}{3-x^2}.
\end{aligned}
\]

\paragraph{在 $0\le x\le1$ 时，偏大偏小一眼看懂}
\[
\begin{array}{c|c|c}
\text{函数}&\text{分式和真正函数值相比}&\text{误差不超过}\\\hline
 e^x&\text{偏小}&3x^5/665\\
 \ln(1+x)&\text{偏小}&x^5/180\\
 \sin x&\text{偏大}&x^5/36\\
 \cos x&\text{偏小}&x^6/288\\
 \tan x&\text{偏小}&x^5/15
\end{array}
\]
\textbf{再记两个有用的结构：}
$E(-x)=1/E(x)$；$L\bigl(-x/(1+x)\bigr)=-L(x)$（$x>-1$）。

\paragraph{做题四步}
\textbf{找 $x$}：对数特别注意 $\ln(1+x)$；
\textbf{看条件}：正切的分母和原函数都要有定义，角度用弧度；
\textbf{算分式}：先求容易计算的估计值；
\textbf{查误差}：确认整个可能范围都能得到同一个舍入结果。

\textbf{一句话收尾：}分式帮你算得快，误差范围帮你算得稳。
\end{summarybox}
